\begin{abstract}
Language models can adapt to new tasks by following instructions provided in context.
Such instructions can be organized as reusable skills, raising the possibility of sharing the internal states they produce across executions.
It remains unclear how the transferability of these states relates to continued dependence on a skill across model layers.
We compare two interventions across layers: transferring hidden states at skill positions from a correct-skill run into an unrelated-skill run, and blocking attention to the skill in the original run.
We also test whether copying skill-position key/value caches, which store information used by attention, can improve transfer.
Across four models and two benchmarks, transfer weakens before dependence on the skill does in all eight settings.
At layer 16, adding caches improves accuracy over state transfer alone by 6.8 to 43.4 percentage points.
Blocking skill reading in subsequent layers reduces the recovered accuracy in each setting, while copying only skill-position caches, even across all layers, falls short of correct-skill accuracy.
These findings distinguish continued skill use from state reuse and show that cached attention information supports partial recovery.
\end{abstract}
